% Sections added only to the VLDB Journal version (make_venues.py inserts them before "Lessons").
% Every number here comes from the same result files as the rest of the paper.
\section{Threats to validity and limitations}\label{sec:limitations}

\paragraph{Hardware and timing.}
All runs use one machine and one search thread. Latency is the per-query median over three rounds
with the method order rotated. Even so, repeated timings of the same searches move by several
percentage points: PercEF's latency saving on DeepImage is $-$8.3\% in the headline run and
$-$13.2\% in the probe-length sweep. We therefore treat latency differences of a few points as
noise, and report distance computations (DC) alongside latency throughout. Multi-threaded
throughput and other processors may rank the methods differently where their margins are small.

\paragraph{Datasets.}
Cohere-1024 and LAION-I2I are subsets that fit in 62\,GB of memory (5 of 10 files and 20 of 31
shards). Our MS MARCO set is a MiniLM-384 stand-in for the embeddings of the Ada-ef paper. The
Big-ANN datasets are 2M-vector prefixes. Results on the full collections may differ, most likely in
absolute savings rather than in which method wins.

\paragraph{Protocol coverage.}
The main protocol fixes $M = 16$, $\mathit{ef}_{\mathrm{construction}} = 500$ and $\tau = 0.95$. We
varied $M$ on four datasets (Table~\ref{tab:vj-m}), the target recall and $k$ on six
(Table~\ref{tab:sweep}), the random split on four (Table~\ref{tab:vj-seeds}), and the probe length on
all 17 same-modality datasets (Table~\ref{tab:vj-probes}). We did not vary
$\mathit{ef}_{\mathrm{construction}}$. Small $k$ remains open: at $k = 10$ neither method saves work.

\paragraph{Baselines.}
DARTH and LAET run in their authors' FAISS fork and are compared with a fixed \ef{} of the same
library, so their numbers are not directly comparable with the HNSWlib ones. DARTH is trained on
1,500 queries, against 10,000 in its paper; more training data could improve its tail but not its
per-step model cost, which is what makes it slow on short searches. PiP is the implementation
shipped with Ada-ef's code; even tuned on our calibration queries, its mean recall on the test
queries reaches the 0.95 target on only 5 of the 20 datasets (Table~\ref{tab:vj-pip}).

\paragraph{Other graph indexes.}
FAISS HNSW and NSG were run through a wrapper around their search call. Neither library can
continue a search after the probe, so every adaptive query pays for a fresh search. We therefore
report only rank correlations for these indexes, not savings. A DiskANN (Vamana) run was not
usable: its returned identifiers were wrong under our software environment. An adaptive search
integrated into these libraries, as it is in HNSWlib, is future work.

\paragraph{Our reimplementation of Ada-ef.}
We run Ada-ef through its authors' code with three changes (\S3.4): a streaming loader for its
statistics, a Boost-free normal quantile checked against SciPy's, and an estimator in the units of
our L2 index. Scores from the streamed and the in-memory statistics are identical on 13 of 17
datasets and differ on some queries, by at most 0.39, on the other four. An earlier version of our
quantile had the wrong sign in the tail (\S\ref{sec:pitfalls}). Every Ada-ef number in this paper
comes from the corrected code.

\section{Guidelines for practitioners}\label{sec:guidelines}

The results suggest a short procedure for choosing a per-query \ef{} method. It needs the raw
vectors, a few hundred held-out queries and a fixed-\ef{} sweep, all of which a deployment has
before it serves traffic.
\begin{enumerate}
\item \textbf{Tune a fixed \ef{} first, and keep it as the baseline.} Compare every adaptive method
  with the best fixed \ef{} at the same mean recall and at the same p1, on held-out queries. Against
  this baseline, adaptive methods saved at most a few percent of the work on most of our datasets,
  and lost it on some.
\item \textbf{Run the KS test on the raw vectors} (seconds). A low KS ($< 0.044$, typical of text
  embeddings) means Ada-ef's model is accurate. A high KS does not by itself mean Ada-ef fails
  (\S\ref{sec:predict}). What it fails on, a calibration that maps to only one or two \ef{} values,
  shows up directly in its table.
\item \textbf{Check the cost of the probe against the length of the search.} Ada-ef adds a
  1025-distance probe and a $d^2$ variance to every query. On searches of several milliseconds with
  $d \le 128$, these are a small share. On high-dimensional data they dominate: even with a
  100-distance probe, Ada-ef is 39\% and 85\% slower than the fixed \ef{} at $d = 1536$ and
  $d = 2048$.
\item \textbf{Choose:}
  \begin{itemize}
  \item Near-Gaussian data, long searches, and tail recall as the priority: \emph{Ada-ef with its
    WAE floor}. It saves 26--27\% of the work on GloVe and is faster than the fixed \ef{} there.
  \item High dimension, short searches, or a predictable cost as the priority: \emph{PercEF}. It is
    cheaper than the tuned fixed \ef{} in 32 of 34 runs and faster in 29, and at its best probe
    length it is faster than the fixed \ef{} on 15 of 17 datasets. A probe of 50--300 distances works
    about equally well.
  \item Queries from another modality than the data, or $k = 10$: \emph{a tuned fixed \ef{}}.
    Neither method saved work there at a cost we would accept.
  \end{itemize}
\item \textbf{Recheck when the data or the queries change.} Both methods' tables are calibrated on
  held-out queries. PercEF's thresholds take seconds to recompute, Ada-ef's statistics a streaming
  pass over the data.
\end{enumerate}
